% TOP_PROBLEM: {"id": 2302074, "title": "Research Problems in Function Theory — Problem 2.74", "category_id": 8, "category": {"id": 8, "name": "analysis", "display_name": "Analysis", "description": "Limits, continuity, calculus, and function theory.", "slug": "analysis", "order_index": 8, "created_at": "2026-07-31T15:26:25.671Z"}, "status": "open", "problem_number": "AMR-022-2074", "set_id": 13, "statusQualification": "The nonpolynomial hypothesis is retained from the motivating statement to exclude an eventually zero Taylor series."}
% FIRST_POSTED: 2026-10-02
% ENTRY_KIND: statement_only
% UnsolvedMath ID 2302074; source catalogue code AMR-022-2074
% !TeX program = lualatex
% Definition and sources only; this file is not an LLM solution attempt.
\documentclass[11pt,a4paper]{article}
\usepackage[margin=25mm]{geometry}
\usepackage{fontspec}
\setmainfont{lmroman10-regular.otf}[BoldFont=lmroman10-bold.otf,
 ItalicFont=lmroman10-italic.otf,BoldItalicFont=lmroman10-bolditalic.otf]
\usepackage{amsmath,amssymb,mathtools}
\usepackage{unicode-math}
\setmathfont{latinmodern-math.otf}
\AtBeginDocument{\let\setminus\smallsetminus}
\usepackage{xurl,hyperref}
\hypersetup{colorlinks=true,urlcolor=blue}
\setcounter{secnumdepth}{0}
\setlength{\parindent}{0pt}
\setlength{\parskip}{5pt}
\setlength{\emergencystretch}{3em}
\begin{document}
\section*{Research Problems in Function Theory — Problem 2.74}\label{2302074}
{\small UnsolvedMath ID 2302074 · AMR-022-2074 · Analysis · AMR Open Problem Lists}
\subsection{Definitions and mathematical statement}
A real entire function has real Taylor coefficients, equivalently $f(\overline z)=\overline{f(z)}$. An entire function has genus at most $q$ when its Hadamard representation uses canonical factors of genus at most $q$ and an exponential of a polynomial of degree at most $q$; finite products are allowed.
For $p\ge0$, let $V_{2p}$ be the real entire functions expressible as
\[
f(z)=e^{-az^{2p+2}}g(z),\qquad a\ge0,
\]
where $g$ is a real entire function of genus at most $2p+1$ whose zeros are all real. Put $U_0=V_0$ and $U_{2p}=V_{2p}\setminus V_{2p-2}$ for $p\ge1$.
Fix $p\ge1$ and a nonpolynomial $f\in U_{2p}$ with $f(0)=1$, and write $f(z)=\sum_{n\ge0}a_nz^n$. Is it true that no $2p+2$ consecutive coefficients vanish? Explicitly, must every $m\ge0$ satisfy
\[
(a_m,a_{m+1},\ldots,a_{m+2p+1})\ne(0,\ldots,0)?
\]
The nonpolynomial condition is necessary because a polynomial has an eventual string of zero coefficients. The question generalizes the two-consecutive-coefficient restriction for the corresponding nonpolynomial class $U_0$.
\subsection{Short English statement}
Does the real-zero class $U_{2p}$ prevent a nonpolynomial entire function from having $2p+2$ consecutive zero Taylor coefficients?
\subsection{Sources}
\begin{itemize}
\item[S1] ULAM.AI and the UnsolvedMath Contributors, supplied problems.json, record 2302074 (AMR-022-2074), AMR Open Problem Lists. Catalogue identity and source transcription; CC BY 4.0.
\url{https://huggingface.co/datasets/ulamai/UnsolvedMath}.
\item[S2] Hayman - Research Problems in Function Theory (2018), Problem 2.74. Original source indexed by AMR-022-2074.
\url{https://arxiv.org/abs/1809.07200}.
\end{itemize}
\end{document}
